\section{Conclusiones y Trabajo Futuro}
\label{sec:conclusiones}

El proyecto tiene implicaciones significativas tanto a nivel técnico como social. Demuestra la viabilidad de utilizar arquitecturas de agentes avanzados en aplicaciones prácticas del mundo real, superando las limitaciones de los chatbots tradicionales. Las herramientas de asesoramiento financiero inteligente pueden ayudar a democratizar el acceso a consejos financieros personalizados, tradicionalmente reservados para personas con altos ingresos o conocimientos especializados \cite{fintech2023}.

En conclusión, el Coach Financiero Generativo representa no solo un éxito técnico en la implementación de agentes inteligentes, sino también una contribución significativa al campo de la tecnología financiera personal, estableciendo un nuevo estándar para las aplicaciones de asesoramiento financiero basadas en IA.

\subsection{Trabajo Futuro}
\label{subsec:trabajo_futuro}

Basándonos en los resultados y limitaciones identificadas, proponemos varias líneas de desarrollo futuro. Una de las más prometedoras es la integración con las prácticas anteriores del curso, creando un ecosistema financiero personal completo. La incorporación de capacidades predictivas mediante los modelos de machine learning desarrollados en prácticas anteriores podría anticipar tendencias de gasto futuras y recomendar estrategias de ahorro proactivas \cite{tool_using_agents2023}.

La extensión del sistema para integrar las herramientas de análisis de datos de la práctica 2 y las capacidades de visualización de la práctica 3 permitiría crear una plataforma unificada. El desarrollo de capacidades multilingües permitiría su aplicación en contextos internacionales, mientras que la integración con APIs bancarias directas podría eliminar la necesidad de importación manual de datos.